\chapter{ DBT (Data Build Tool): The Modern Way to Transform Data}
\begin{mybox}{Anggota kelompok}
\begin{enumerate}
    \itemsep0em
	\item ENGELBERTUS RANDE (NIM 25/573149/PPA/07211)
	\item GILBERT FANDILIAM MOOY (NIM 25/574767/PPA/07259)
	\item Sinta Siti Nuriah (NIM 25/575177/PPA/07274)
	\item Vicky Mahfudy (NIM 25/573019/PPA/07207)
\end{enumerate}
\end{mybox}

\section{Pengenalan DBT}
\subsection{Konsep dan Arsitektur DBT }
\textbf{DBT } (\textit{data build tool}) adalah \textit{command-line tool }dan \textit{framework }berbasis Python yang memungkinkan tim data melakukan transformasi data di dalam \textit{data warehouse }menggunakan SQL murni (1). DBT bekerja pada lapisan \textbf{T } (\textit{Transform}) dalam paradigma \textbf{ELT } (\textit{Extract, }\textit{Load, Transform}).
\begin{mybox}{Definisi Resmi DBT Labs}
\textit{``dbt enables data analysts and engineers to transform data
in their warehouses more effectively. ''} --- dbt Labs, 2024.
 
Saat ini lebih dari \textbf{40.000 perusahaan} menggunakan DBT dalam
lingkungan produksi, menjadikannya standar industri \textit{de facto}
untuk transformasi data di atas \textit{cloud data warehouse}.
\end{mybox}
\subsubsection{Paradigma ELT dan Posisi DBT}
 
Berbeda dengan \textbf{ETL} tradisional yang mentransformasi data
\textit{sebelum} memuatnya ke \textit{warehouse}, paradigma
\textbf{ELT} modern memuat data mentah terlebih dahulu lalu
mentransformasinya \textit{di dalam} warehouse menggunakan
kekuatan komputasi cloud.

\begin{center}
\begin{tikzpicture}[
  box/.style={rectangle,draw=ugmblue,fill=blue!8,rounded corners=5pt,
    minimum width=2.5cm,minimum height=1.1cm,align=center,font=\small\bfseries},
  act/.style={rectangle,draw=ugmblue,fill=ugmblue,rounded corners=5pt,
    minimum width=2.5cm,minimum height=1.1cm,align=center,
    font=\small\bfseries,text=white},
  arr/.style={->,thick,color=ugmgold,line width=1.5pt}
]
  \node[box,fill=red!10,draw=red!50] (E)
    {Extract\\{\tiny Fivetran/Airbyte}};
  \node[box,fill=orange!10,draw=orange!50,right=1.2cm of E] (L)
    {Load\\{\tiny Raw tables}};
  \node[act,right=1.2cm of L] (T)
    {\textcolor{ugmgold}{T}ransform\\{\tiny \textbf{DBT} bekerja di sini}};
  \node[box,fill=green!10,draw=green!50,right=1.2cm of T] (A)
    {Analyse\\{\tiny BI / Dashboard}};
  \draw[arr] (E)--(L);
  \draw[arr] (L)--(T);
  \draw[arr] (T)--(A);
\end{tikzpicture}
\end{center}
DBT tidak mengekstrak atau memuat data---tugas tersebut dilakukan
oleh \textit{ingestion tool} seperti Fivetran atau Airbyte. DBT
hanya bekerja pada data yang sudah ada di dalam \textit{warehouse}
\citep{bi4all2024}. Pendekatan ini melahirkan peran baru bernama
\textbf{\textit{analytics engineer}}: menggabungkan keahlian SQL
seorang \textit{data analyst} dengan praktik rekayasa perangkat
lunak.
\subsubsection{Arsitektur Berlapis (\textit{Layered Architecture})}
 
Praktik terbaik dalam DBT menggunakan arsitektur berlapis yang
memisahkan tanggung jawab setiap tahap transformasi
\citep{dbtdoc}:
 
\begin{center}
\begin{tikzpicture}[
  layer/.style={rectangle,draw=ugmblue!40,rounded corners=4pt,
    minimum width=11cm,minimum height=0.95cm,
    align=center,font=\small\bfseries},
  arr/.style={->,thick,color=ugmgold}
]
  \node[layer,fill=red!10] (raw) at (0,0)
    {RAW / SOURCES \quad
     {\normalfont\small Data mentah dari sistem sumber, tidak disentuh DBT}};
  \node[layer,fill=orange!10] (stg) at (0,-1.35)
    {STAGING \texttt{(stg\_*)} \quad
     {\normalfont\small Pembersihan, rename kolom, standardisasi tipe data}};
  \node[layer,fill=yellow!25] (int) at (0,-2.7)
    {INTERMEDIATE \texttt{(int\_*)} \quad
     {\normalfont\small Logika bisnis menengah: join, kalkulasi, flag}};
  \node[layer,fill=green!10] (mrt) at (0,-4.05)
    {MARTS \texttt{(fct\_* / dim\_*)} \quad
     {\normalfont\small Tabel analitik final: fact tables dan dimension tables}};
  \draw[arr] (raw)--(stg);
  \draw[arr] (stg)--(int);
  \draw[arr] (int)--(mrt);
\end{tikzpicture}
\end{center}

\section{Model dalam DBT dan Pemanfaatannya}
\subsection{Classic Model}
Pada modul ini digunakan database \textbf{ClassicModels} dari
MySQL Tutorial \citep{mysqltut} sebagai dataset eksplorasi.
ClassicModels adalah database peritel miniatur mobil klasik
yang menyediakan data bisnis nyata dengan relasi antar tabel
yang representatif untuk latihan membangun \textit{data warehouse}.

\begin{table}[H]
\centering
\caption{Delapan Tabel Sumber ClassicModels (OLTP)}
\begin{tabularx}{\linewidth}{llX}
\toprule
\textbf{Tabel} & \textbf{Baris} & \textbf{Deskripsi} \\
\midrule
\texttt{customers}    & 122   & Pelanggan: nama, kota, negara, \textit{credit limit} \\
\texttt{orders}       & 326   & Header pesanan: tanggal, status \\
\texttt{orderdetails} & 2.996 & Item pesanan: produk, kuantitas, harga satuan \\
\texttt{products}     & 110   & Katalog produk: nama, lini, vendor, harga \\
\texttt{productlines} & 7     & Lini produk: Classic Cars, Motorcycles, dll. \\
\texttt{employees}    & 23    & Karyawan dan struktur organisasi \\
\texttt{offices}      & 7     & Kantor penjualan: kota, negara, wilayah \\
\texttt{payments}     & 273   & Pembayaran pelanggan \\
\bottomrule
\end{tabularx}
\end{table}

Dari 8 tabel OLTP ini akan dibangun \textbf{star schema} dengan 5 tabel dimensi dan 2 tabel fakta menggunakan DBT\@.
\subsection{Struktur Model dalam DBT}
Dalam DBT, \textit{best practices} perancangan struktur model menggunakan arsitektur berlapis (\textit{layered architecture}) untuk memisahkan tanggung jawab dari setiap tahap transformasi. Pada implementasi proyek \texttt{classicmodels\_dbt}, strukturnya dibagi sebagai berikut:

% \subsection*{A. Direktori Project (\texttt{models/})}
DBT menyusun \textit{pipeline}-nya dalam folder \texttt{models/} yang mencakup dua subfolder utama, yaitu:

\begin{itemize}
    \item \textbf{Folder \texttt{staging/} (Layer Staging)} \\
    Layer ini bertujuan untuk melakukan standardisasi kolom, perbaikan tipe data, dan pembersihan awal dari \textit{raw data}. Model ini di-\textit{materialize} sebagai \textit{view}.
    \begin{itemize}
        \item Terdapat file konfigurasi \texttt{schema.yml} (untuk mendefinisikan sumber data dari OLTP asli beserta \textit{testing} dasar seperti \texttt{not\_null} atau \texttt{unique}).
        \item File-file SQL yang diberi prefiks \texttt{stg\_} (misalnya: \texttt{stg\_customers.sql}, \\
        \texttt{stg\_orders.sql}, \texttt{stg\_orderdetails.sql}, \texttt{stg\_products.sql}, \\
        \texttt{stg\_employees.sql}, \texttt{stg\_offices.sql}).
    \end{itemize}

    \item \textbf{Folder \texttt{marts/} (Layer Marts / Core)} \\
    Layer ini memuat struktur tabel analitik final yang menggabungkan logika bisnis dan menerapkan desain \textit{Star Schema} (\textit{Dimensional Modeling}). Tabel-tabel di sini di-\textit{materialize} sebagai tabel fisik di dalam database.
    \begin{itemize}
        \item File berprefiks \texttt{dim\_} (Tabel Dimensi / Konteks): \texttt{dim\_date.sql}, \\
        \texttt{dim\_product.sql}, \texttt{dim\_customer.sql}, \texttt{dim\_employee.sql}, \\
        \texttt{dim\_office.sql}.
        \item File berprefiks \texttt{fct\_} (Tabel Fakta / Pengukuran Bisnis): \\
        \texttt{fct\_sales.sql}, \texttt{fct\_monthly\_sales.sql}.
        \item File \texttt{schema.yml} untuk mengatur dokumentasi dan aturan \textit{testing} di level \textit{marts}.
    \end{itemize}
\end{itemize}

% \subsection*{B. Folder Pendukung Lainnya}
Selain folder \texttt{models/}, proyek ini juga memanfaatkan struktur komponen DBT lainnya:

\begin{itemize}
    \item \texttt{tests/}: Berisi \textit{custom SQL script} untuk pengujian logika bisnis spesifik (contohnya memastikan \textit{revenue} selalu positif, \texttt{assert\_revenue\_positive.sql}).
    \item \texttt{macros/}: Berisi fungsi \textit{reusable} dalam bahasa Jinja2 yang dipakai berulang dalam kode SQL (misalnya \texttt{revenue\_segment.sql}).
    \item \texttt{dbt\_project.yml}: Merupakan file konfigurasi induk yang mendefinisikan \textit{paths} (alur folder proyek) serta pengaturan \textit{materialization} di tiap layer.
\end{itemize}

\section{Data Quality dalam DBT}
% \section{Pengendalian Kualitas Data (\textit{Data Quality}) pada DBT}

Dalam perancangan \textit{Data Warehouse}, memastikan keakuratan dan integritas data adalah sebuah keharusan. DBT (\textit{Data Build Tool}) memfasilitasi kebutuhan ini melalui fitur \textbf{Pengujian Otomatis (\textit{Automated Testing})} yang terintegrasi secara mulus di dalam \textit{pipeline} transformasi data.

Konsep dasar pengujian di dalam DBT sangatlah sederhana: \textbf{sebuah tes adalah \textit{query} SQL (\texttt{SELECT}) yang secara khusus dirancang untuk menangkap data anomali atau data yang melanggar aturan}. Apabila \textit{query} tersebut tidak menemukan data yang salah (mengembalikan 0 baris), maka tes tersebut dinyatakan \textbf{lulus (\textit{pass})}. 

Secara umum, DBT membagi pengujian ke dalam dua pendekatan utama:

\subsection{Pengujian Bawaan (\textit{Generic Tests})}
\textit{Generic Tests} adalah sekumpulan fungsi pengujian standar bawaan DBT\@. Pengujian ini sangat praktis karena tidak perlu menulis kode SQL secara manual; cukup dengan mendeklarasikannya di dalam file konfigurasi (seperti \texttt{schema.yml})\@. 

Terdapat empat \textit{generic tests} utama yang sering digunakan di industri:
\begin{itemize}
    \item \textbf{\texttt{not\_null}}: Memvalidasi bahwa sebuah kolom wajib memiliki isi (tidak boleh \textit{null}).
    \item \textbf{\texttt{unique}}: Menjamin tidak ada nilai ganda (duplikat) di dalam sebuah kolom. Aturan ini sangat penting untuk memvalidasi kolom \textit{Primary Key}.
    \item \textbf{\texttt{accepted\_values}}: Membatasi nilai data agar hanya sesuai dengan daftar validasi yang telah ditetapkan (misalnya, kolom \textit{status} hanya boleh berisi nilai `Shipped', `Resolved', atau `Cancelled').
    \item \textbf{\texttt{relationships}}: Memverifikasi integritas referensial antar tabel (identik dengan konsep \textit{Foreign Key} pada \textit{relational database}).
\end{itemize}

\textbf{Contoh Pendefinisian di \texttt{schema.yml}:}
\begin{lstlisting}[caption=Konfigurasi Generic Test]
sources:
  - name: raw
    database: classicmodels
    tables:
      - name: orders
        columns:
          - name: orderNumber
            tests: 
              - not_null
              - unique
\end{lstlisting}

\subsection{Pengujian Logika Bisnis (\textit{Singular Tests})}
Dalam skenario dunia nyata, aturan bisnis seringkali lebih kompleks dari sekadar mengecek nilai kosong atau duplikat. Untuk mengatasi hal ini, DBT menyediakan \textit{Singular Tests}. 

Pada pendekatan ini, pengembang dapat menulis skrip SQL (\textit{custom query}) yang spesifik dan menyimpannya di dalam direktori \texttt{tests/}. Sebagai contoh, pada studi kasus \textbf{Classic Models}, kita harus memastikan bahwa perhitungan pendapatan (\textit{revenue}) tidak boleh bernilai negatif. 

\textbf{Contoh Skrip \texttt{tests/assert\_revenue\_positive.sql}:}
\begin{lstlisting}[caption=Skrip Singular Test]
-- Skrip ini dirancang untuk mendeteksi baris dengan pendapatan negatif.
-- Tes akan GAGAL jika ada output yang dihasilkan.
SELECT
    order_number,
    line_revenue
FROM {{ ref('fct_sales') }}
WHERE line_revenue < 0
\end{lstlisting}

\subsection{Menjalankan Pengujian (Eksekusi)}
Setelah semua aturan pengujian (\textit{generic} maupun \textit{singular}) didefinisikan, langkah terakhir adalah mengeksekusinya melalui antarmuka \textit{Command-Line} (CLI):
\begin{itemize}
    \item \texttt{dbt test}: Untuk menjalankan seluruh skrip pengujian saja.
    \item \texttt{dbt build}: Untuk menjalankan proses pembuatan tabel (\textit{run}) dan pengujian (\textit{test}) sekaligus dalam satu urutan kerja.
\end{itemize}

% \subsection*{Studi Kasus: Classic Models}
% Pada proyek pembelajaran Classic Models, implementasi arsitektur ini terbukti sangat efektif. \textit{Pipeline} yang dibangun berhasil memvalidasi seluruh logika dengan status \textbf{18/18 tes lulus (\textit{passed})}. Hal ini menjadi jaminan bahwa data telah bertransformasi dari sistem sumber (OLTP) ke dalam bentuk \textit{Star Schema} (\textit{Marts}) dengan tingkat akurasi dan kualitas yang tinggi.